\chapter{研究与实现路线}

\section{两条并行轨道}

Nous Wave 的长期目标继续按完整 Subject cognition/personality system 研究，不因为中期检查而削弱；工程实施则选择边界最成熟的 owner 做纵向闭环。当前形成两条同步但不同节奏的轨道：

\begin{description}
\item[长期设计轨道] 继续收敛 Self constitution、Motivation、Social Spec、Episode/Journal/Dream、检索算法与 Heptalogos 集成，保持从宏观 ownership 到 executable contract 的完整推导。
\item[中期工程轨道] 在 2026 年 10 月 27 日前实现 Memory Reference Profile，从 Observation/Evidence 一直到持久化、检索、Use feedback、生命周期、恢复和演示形成可验证闭环。
\end{description}

这不是 MVP 式删除未来能力。Reference Profile 只决定当前施工顺序；未进入中期实现的 Cognition domain 仍保留其目标语义，不建立空 crate、空表或虚假“已支持”接口。

\section{Memory Reference Profile 的依赖顺序}

\begin{enumerate}
\item \textbf{重写当前 Architecture/Decision 基线：} 用本报告中的 Memory、Schema、Social 和 algorithm status 替换仓库上一代 MemoryClass/Anchor/默认算法叙述；明确哪些当前代码资产直接继承。
\item \textbf{冻结 Memory semantic contract：} MemoryObject/Revision、\texttt{cognitive\_role}、\texttt{formation\_mode}、Evidence/aboutness、time、revision/lifecycle、AccessPolicy、Use semantics、suppression/purge。
\item \textbf{冻结 Retrieval contract：} QueryPlan lanes/per-lane budgets、multi-view object collapse、weighted RRF、hard filters、final Authority revalidation、diagnostics/degradation。
\item \textbf{冻结可靠写入合同：} request/event identity、CAS、idempotency、transaction ordering、authority sequence/epoch、projection invalidation/outbox、水位和 timeout-after-commit。
\item \textbf{下降为 executable Spec：} SQL DDL、Proto、Rust domain types、TypeScript Client/Host interfaces、error codes、module ownership、acceptance scenarios。
\item \textbf{实现/重构：} 保留已经正确的 Evidence/Revision/Host-Kernel/Serving infrastructure，替换旧 MemoryClass、Anchor Authority、candidate-dependent ranking、non-idempotent ReportUse 和 hard AUTO gating。
\item \textbf{资格与演示：} 真实 PostgreSQL 端到端形成、修订/矛盾、混合检索、Use feedback、restart、suppress/restore/purge 和 provenance tracing；明确 PASS/FAIL/NOT\_RUN 范围。
\end{enumerate}

\section{中期非阻塞实验}

EPA、Residual、Wave/PPR、advanced accessibility、dense backend benchmark 可以并行推进，但任何单项失败都不能阻塞 Memory Authority/Revision/Retrieval baseline 的正确性。Memory Reference Profile 首先依赖 Exact/Entity/Lexical/Dense 的强基线。

若核心闭环提前完成，可以增加 Episode 或一个小型 Social prototype 作为演示扩展；不能反过来为了展示面积牺牲 Memory 的事务、来源和可靠性合同。

\section{中期之后的实现顺序}

在不改变整体设计研究优先级的前提下，当前较合理的工程顺序为：

\begin{enumerate}
\item Memory Reference Profile；
\item Social Cognition：RelationshipGraph；
\item Social Cognition：LanguageConvention；
\item Motivation；
\item Self/Persona；
\item Episode/Journal/CognitiveSchema 的完整高级维护与 Dream/cross-domain cognition。
\end{enumerate}

Social 的首批对象大量复用 SubjectId、EntityRef、Evidence、Revision、EpistemicClass、valid time、proposal 和 purge 等 Memory/Cognition Foundation，因此适合作为 Memory 之后验证跨 owner 复用能力的第二个领域。Self constitution 和跨领域 Dream 对 Authority 要求更高，后实施不代表其研究价值更低。

\section{Heptalogos 集成顺序}

跨库集成不等待所有高级 cognition 完成。Memory Reference Profile 稳定后即可建立最小真实链路：

\begin{center}
H canonical interaction fact $\rightarrow$ NW Observation\\
NW query/context $\rightarrow$ H Reaction/Activity\\
H use/outcome $\rightarrow$ NW idempotent receipt
\end{center}

随后再增加 Motivation Goal $\leftrightarrow$ H Commitment 和 maintenance lease/offline cognition。H Foundation 已经保留高级 cognition contribution seam；下一步重点是把 Nous Wave 确立为这些 cognition capability 的具体 owner/provider，而不是在 H 内重复实现一套 Memory/Social/Self ontology。

\section{算法采用条件}

候选算法只有在长期质量--成本比较中稳定优于强基线后，才考虑成为默认。算法是否存在实现、是否和某个外部项目同构、或是否在合成数据中保持数学性质，都不等于真实 retrieval utility 已经成立。所有默认参数记录适用 workload 和证据边界。\cite{nous-research-package}

\section{阶段工作的记录方式}

每项工作记录：研究问题、语义 owner、来源/版本、候选方案、实验/实现路径、执行状态、结果、限制及是否进入当前 Spec。实现代码、设计结论、实验观察、运行资格和外部研究分别标注，避免“代码已存在”被误写成“方案已验证”。